% GENERATED FILE. Edit canonical lessons, then run npm run generate:books.
% canonical-source-hash: fnv1a64:3fcbf13930aae210
% canonical-lessons: SA-C174-corayati, SA-C174-bhaksayati, SA-C174-tadayati, SA-C174-palayati, SA-C174-pidayati

\chapter{Steals, Eats up, Strikes, Keeps, Hurts}
\label{ch:sa-steals-eats-up-strikes-keeps-hurts}
\hlchaptermodality{174}

\begin{chapteropening}
\textbf{By the end of this chapter:} \emph{I can say steals, eats up, strikes, keeps and hurts.}

Complete the last lesson of chapter 174: I can say steals, eats up, strikes, keeps and hurts.
\end{chapteropening}

\section[\texorpdfstring{\sk{चोरयति} (corayati)}{चोरयति (corayati)}]{\sk{चोरयति} (corayati) — steals}

\label{lesson:SA-C174-corayati}



\begin{quote}
Before the new one: say the Sanskrit for serves, then the Sanskrit for enjoys.
\end{quote}

\subsection*{You'll want to know: \sk{चोरयति}}
\textbf{\sk{चोरयति}} — \emph{corayati} — \textquotedblleft{}steals\textquotedblright{}.

\begin{grammarlens}[title={What you've built}]
One more word for this chapter.
\end{grammarlens}

\subsection*{Your turn}
\begin{itemize}
\raggedright
  \item \emph{Say it:} \emph{corayati}
  \item \emph{Say it:} \emph{corayati}, once more
  \item \emph{Say it:} say \emph{corayati}
  \item \emph{Recall:} say the Sanskrit for serves, then the Sanskrit for enjoys, then say \emph{corayati} again
\end{itemize}

\begin{tcolorbox}[breakable,colback=teal!4,colframe=teal!35!black,title={Before you move on}]
What does \sk{चोरयति} mean? (\textquotedblleft{}Steals\textquotedblright{}.) Say it once more.
\end{tcolorbox}
\section[\texorpdfstring{\sk{भक्षयति} (bhakṣayati)}{भक्षयति (bhakṣayati)}]{\sk{भक्षयति} (bhakṣayati) — eats up, devours}

\label{lesson:SA-C174-bhaksayati}



\begin{quote}
Before the new one: say the Sanskrit for enjoys, then the Sanskrit for steals.
\end{quote}

\subsection*{You'll want to know: \sk{भक्षयति}}
\textbf{\sk{भक्षयति}} — \emph{bhakṣayati} — \textquotedblleft{}eats up, devours\textquotedblright{}.

\begin{grammarlens}[title={What you've built}]
One more word for this chapter.
\end{grammarlens}

\subsection*{Your turn}
\begin{itemize}
\raggedright
  \item \emph{Say it:} \emph{bhakṣayati}
  \item \emph{Say it:} \emph{bhakṣayati}, once more
  \item \emph{Say it:} say \emph{bhakṣayati}
  \item \emph{Recall:} say the Sanskrit for enjoys, then the Sanskrit for steals, then say \emph{bhakṣayati} again
\end{itemize}

\begin{tcolorbox}[breakable,colback=teal!4,colframe=teal!35!black,title={Before you move on}]
What does \sk{भक्षयति} mean? (\textquotedblleft{}Eats up, devours\textquotedblright{}.) Say it once more.
\end{tcolorbox}
\section[\texorpdfstring{\sk{ताडयति} (tāḍayati)}{ताडयति (tāḍayati)}]{\sk{ताडयति} (tāḍayati) — strikes, beats}

\label{lesson:SA-C174-tadayati}



\begin{quote}
Before the new one: say the Sanskrit for steals, then the Sanskrit for eats up.
\end{quote}

\subsection*{You'll want to know: \sk{ताडयति}}
\textbf{\sk{ताडयति}} — \emph{tāḍayati} — \textquotedblleft{}strikes, beats\textquotedblright{}.

\begin{grammarlens}[title={What you've built}]
One more word for this chapter.
\end{grammarlens}

\subsection*{Your turn}
\begin{itemize}
\raggedright
  \item \emph{Say it:} \emph{tāḍayati}
  \item \emph{Say it:} \emph{tāḍayati}, once more
  \item \emph{Say it:} say \emph{tāḍayati}
  \item \emph{Recall:} say the Sanskrit for steals, then the Sanskrit for eats up, then say \emph{tāḍayati} again
\end{itemize}

\begin{tcolorbox}[breakable,colback=teal!4,colframe=teal!35!black,title={Before you move on}]
What does \sk{ताडयति} mean? (\textquotedblleft{}Strikes, beats\textquotedblright{}.) Say it once more.
\end{tcolorbox}
\section[\texorpdfstring{\sk{पालयति} (pālayati)}{पालयति (pālayati)}]{\sk{पालयति} (pālayati) — keeps, protects}

\label{lesson:SA-C174-palayati}



\begin{quote}
Before the new one: say the Sanskrit for eats up, then the Sanskrit for strikes.
\end{quote}

\subsection*{You'll want to know: \sk{पालयति}}
\textbf{\sk{पालयति}} — \emph{pālayati} — \textquotedblleft{}keeps, protects\textquotedblright{}.

\begin{grammarlens}[title={What you've built}]
One more word for this chapter.
\end{grammarlens}

\subsection*{Your turn}
\begin{itemize}
\raggedright
  \item \emph{Say it:} \emph{pālayati}
  \item \emph{Say it:} \emph{pālayati}, once more
  \item \emph{Say it:} say \emph{pālayati}
  \item \emph{Recall:} say the Sanskrit for eats up, then the Sanskrit for strikes, then say \emph{pālayati} again
\end{itemize}

\begin{tcolorbox}[breakable,colback=teal!4,colframe=teal!35!black,title={Before you move on}]
What does \sk{पालयति} mean? (\textquotedblleft{}Keeps, protects\textquotedblright{}.) Say it once more.
\end{tcolorbox}
\section[\texorpdfstring{\sk{पीडयति} (pīḍayati)}{पीडयति (pīḍayati)}]{\sk{पीडयति} (pīḍayati) — hurts, troubles}

\label{lesson:SA-C174-pidayati}



\begin{quote}
Before the new one: say the Sanskrit for strikes, then the Sanskrit for keeps.
\end{quote}

\subsection*{You'll want to know: \sk{पीडयति}}
\textbf{\sk{पीडयति}} — \emph{pīḍayati} — \textquotedblleft{}hurts, troubles\textquotedblright{}.

\begin{grammarlens}[title={What you've built}]
One more word for this chapter.
\end{grammarlens}

\subsection*{Your turn}
\begin{itemize}
\raggedright
  \item \emph{Say it:} \emph{pīḍayati}
  \item \emph{Say it:} \emph{pīḍayati}, once more
  \item \emph{Say it:} say \emph{pīḍayati}
  \item \emph{Recall:} say the Sanskrit for strikes, then the Sanskrit for keeps, then say \emph{pīḍayati} again
\end{itemize}

\begin{tcolorbox}[breakable,colback=teal!4,colframe=teal!35!black,title={Before you move on}]
What does \sk{पीडयति} mean? (\textquotedblleft{}Hurts, troubles\textquotedblright{}.) Say it once more.
\end{tcolorbox}
